\section*{Chapitre \ref{ch:ml:validation} --- Validation}

\begin{solution}{exo:ml:validation:1}
$0.30~\% + 1.16\times0.10~\% = 0.42~\%$~: le score du gagnant surestime de 0.12 point le vrai 0.30~\% de chaque candidat.
\end{solution}

\begin{solution}{exo:ml:validation:2}
Les étiquettes du test couvrent les mois 100 à 141. La purge supprime les étiquettes d'entraînement commençant aux mois 98 et 99 (elles se terminent
à l'intérieur de la période) et aux mois 140 et 141 (elles commencent à l'intérieur de celle-ci)~; l'embargo supprime aussi le mois 142.
\end{solution}

\begin{solution}{exo:ml:validation:3}
Moyenne 0.18~\%, écart-type 0.03~\%, erreur type $0.03/\sqrt3 = 0.017~\%$~: dix erreurs types au-dessus de zéro sur
des cibles impossibles à prédire. Le pipeline fuit.
\end{solution}

\begin{solution}{exo:ml:validation:4}
La jointure à la fin de période fait fuir de l'information sur les rendements réels (le rapport les déplace)~; sur une cible de bruit, elle n'apporte
rien, si bien que le canari ne voit aucun gain ($-0.44~\%$). Le test de troncature demande à quel moment chaque valeur est devenue connue et repère la
surprise un mois trop tôt. L'encodage de la cible est construit à partir de celle-ci~: sur une cible de bruit, il fuit encore (le score du
canari, 4.03~\%), et recalculé sur des données tronquées à une date, il diffère (le test de troncature).
\end{solution}

\begin{solution}{exo:ml:validation:5}
Les variables situent la période, si bien que des plis mélangés fuiraient (les voisins sont reconnaissables)~; et la distribution des
variables sur la période de test diffère de celle de l'entraînement (dérive, nouveau régime, source de données modifiée). Cela ne montre pas que le
modèle se trompe ni que les variables prédisent les rendements~: une variable peut se déplacer tout en restant utile.
\end{solution}

\begin{solution}{exo:ml:validation:6}
Les statistiques calculées sur tout l'échantillon utilisent les données de test~: l'échelle et le centre de chaque variable dans les lignes d'entraînement dépendent du
futur. Pour des variables stationnaires et bien conditionnées, la fuite est faible~; pour des variables tendancielles ou dérivantes (volumes, prix,
comptages), elle ne l'est pas, et c'est un biais d'anticipation, quelle que soit sa taille. Ajuster le normalisateur dans chaque pli d'entraînement, ou utiliser
des rangs en coupe transversale calculés date par date.
\end{solution}

\begin{solution}{exo:ml:validation:7}
\texttt{announcement\_effect(0.04)}~: la jointure à la fin de période obtient 5.47~\% (1.27~\% pour 2~\%), la jointure à la date de dépôt $-0.14\%$.
La variance expliquée par une surprise connue est proportionnelle au carré de l'effet, si bien que le doubler multiplie
le gain par environ quatre~; l'erreur propre du modèle ajoute le même petit terme négatif dans les deux cas.
\end{solution}

\begin{solution}{exo:ml:validation:8}
Dans chaque pli externe, la procédure (recherche, choix, réajustement) n'utilise que les lignes d'entraînement externes, si bien que le modèle ajusté
est indépendant des lignes de test externes~; son score y est une estimation sans biais du score de généralisation de ce que la
procédure a produit à partir de ce volume de données, et il en va de même pour la moyenne sur les plis. Le meilleur à plat prend un maximum sur les candidats
de scores mesurés sur les mêmes lignes de test, et d'après la \cref{prop:ml:validation:max}, son espérance dépasse le
score vrai de chaque candidat.
\end{solution}

\begin{solution}{pb:ml:validation:1}
\begin{enumerate}
\item $-0.19~\%$~: rien n'est prévisible à l'instant de décision, si bien que le meilleur score possible est zéro et qu'un modèle ajusté ne peut que
faire moins bien que lui.
\item Le rendement de la dernière ligne d'entraînement est celui du mois 160~: un écart empêche les étiquettes d'entraînement de se terminer à l'intérieur du test.
\item Date de dépôt~: la dernière surprise déposée à la fin du mois $t$ (trimestre clos en $t-1$ ou avant). Fin de période~: la
surprise du trimestre clos en $t$, déposée pendant $t+1$.
\item 2~\% de rendement par unité de surprise standardisée.
\item 8.97~\%~: le rendement sectoriel du même mois explique la part du facteur sectoriel dans la variance, qui est grande.
\item 0.13~\% sur l'ensemble de l'échantillon, $-0.16~\%$ sur les mois d'entraînement.
\item 38\,080 contre aucun.
\item 3.47~\% et $-3.37~\%$~; des attributs persistants permettent aux arbres profonds de reconnaître un titre sur une période donnée et
de lire les étiquettes chevauchantes de ses voisins.
\item De $-0.60~\%$ à $-0.27~\%$~; le meilleur est 0.13 point au-dessus de la médiane $-0.40~\%$.
\item $-0.20~\%$, contre une médiane de $-0.19~\%$~; les scores des deux périodes sont corrélés à $-0.16$.
\item Meilleur à plat $-0.62~\%$, imbriqué $-0.85~\%$.
\item Le maximum d'estimations bruitées surestime en espérance~; le score imbriqué est mesuré sur des données que le choix n'a jamais
vues.
\item Jointure à la fin de période~: test de troncature (4.77 contre 0). Encodage de la cible~: canari (4.03~\%) et test de troncature.
Criblage~: canari (0.17~\%) et test de troncature (5.13). Chevauchement~: contrôle de chevauchement des plis (38\,080 contre 0). Sélection~:
échantillon de test intact ou recherche imbriquée.
\item Chaque détecteur regarde un seul endroit (variables, pipeline, découpages, choix)~; une fuite réside en un seul endroit.
\item Les attributs identifient la période, si bien que toute découpe mélangée fuit et que les mois de test diffèrent des mois
d'entraînement.
\item \emph{Résultat nommé}~: sur une cible impossible à prédire, la jointure à la fin de période affiche un gain fictif de 1.27~\%, l'encodage de la cible 8.97~\%,
le criblage sur tout l'échantillon 0.13~\% ($-0.16~\%$ honnête), les plis mélangés sur des étiquettes chevauchantes 3.47~\% ($-3.37~\%$ après purge), et
la sélection du meilleur graine sur le test 0.13 point~; la troncature, le canari, le contrôle de chevauchement et l'échantillon intact les repèrent dans cet ordre.
\item Les variables recalculées aux dates de l'échantillon coïncident~; le score du canari est à moins de deux erreurs types du hasard~; aucune
étiquette d'entraînement ne chevauche une étiquette de test~; l'AUC de la \omterm{def:ml:validation:audit}{validation adverse} est consigné~; l'échantillon de test est resté intact~; le nombre d'essais est consigné.
\item Aucune ne survit~: chacune disparaît dès que les données arrivent en temps réel. La sélection se traduit par un score en conditions réelles plus faible~;
les autres font perdre le trading réel d'emblée, ce qui explique qu'on les découvre tard si le trading simulé (paper trading) est court.
\item Un constat accompagné de son effet mesuré, de sa correction et de chaque résultat qu'il invalide.
\item Estimer honnêtement les performances d'une procédure sur des données qu'elle n'a pas encore vues, ce qui exige qu'elle ne les ait réellement pas
vues.
\end{enumerate}
\end{solution}

\begin{solution}{iq:ml:validation:1}
Information présente dans l'entraînement ou dans les variables, et qui ne serait pas disponible à l'instant de décision. Exemples~: fondamentaux joints à la fin de
période au lieu de la date de dépôt~; données macroéconomiques révisées utilisées telles qu'initialement publiées~; univers des membres actuels de l'indice
(biais de survie)~; variables normalisées avec des statistiques futures~; étiquettes chevauchantes entre plis mélangés.
\iqlookfor{une définition précise liée à l'instant de décision, et des cas financiers concrets.}
\end{solution}

\begin{solution}{iq:ml:validation:2}
Le meilleur des vingt moyennes sur les plis est biaisé à la hausse (maximum d'estimations bruitées). Présenter une \omterm{def:ml:validation:nested}{validation croisée imbriquée} ou une
analyse glissante (walk-forward) de la procédure entière, ainsi que le nombre de réglages testés.
\iqlookfor{biais de sélection et remède par validation imbriquée ou échantillon de test.}
\end{solution}

\begin{solution}{iq:ml:validation:3}
Recalculer chaque variable sur des données tronquées à plusieurs dates et comparer avec la variable stockée (toute différence signale
un biais d'anticipation)~; lancer le pipeline sur des cibles de bruit de même structure (canari)~; vérifier les temps de connaissance des entrées
par rapport aux instants de décision.
\iqlookfor{tests en boîte noire qui ne reposent pas sur la lecture du code.}
\end{solution}

\begin{solution}{iq:ml:validation:4}
Un classifieur entraîné à séparer les lignes d'entraînement des lignes de test~; son AUC mesure à quel point elles diffèrent. Une AUC élevée signale
que les deux côtés du découpage diffèrent (dérive) et que les variables situent le temps, si bien que les découpages mélangés fuient~; les variables sur lesquelles il s'appuie
sont celles qu'il faut examiner.
\iqlookfor{ce qu'elle mesure et ses deux usages.}
\end{solution}

\begin{solution}{iq:ml:validation:5}
Les encodages calculés avec des cibles de la même période ou de périodes futures fuient la cible elle-même. Les calculer à partir des seules périodes passées
(moyennes expansives ou décalées), à l'intérieur de chaque pli d'entraînement, avec un lissage vers la moyenne globale.
\iqlookfor{décaler la statistique, et la calculer à l'intérieur du pli.}
\end{solution}

\begin{solution}{iq:ml:validation:6}
Lancer l'\omterm{def:ml:validation:audit}{audit de fuite} (troncature, canari, chevauchement des plis)~: une fuite s'y révèle. Vérifier ensuite la recherche~: nombre d'essais,
score imbriqué, statistiques déflatées~; un score imbriqué proche de zéro signifie \omterm{def:ml:why-financial-machine-learning-is-different:generalisation}{surapprentissage}. Si les deux passent, comparer les distributions des variables
et les erreurs du modèle en conditions réelles avec le backtest (\omterm{def:ml:validation:audit}{validation adverse}, tests de dérive, chapitre~12), puis calculer
combien de temps doit durer la période de trading simulé (paper trading) pour que l'écart soit significatif.
\iqlookfor{un diagnostic ordonné avec un test pour chaque cause.}
\end{solution}
